
\documentclass{article}
\usepackage{colortbl}
\usepackage{makecell}
\usepackage{multirow}
\usepackage{supertabular}

\begin{document}

\newcounter{utterance}

\centering \large Interaction Transcript for game `hot\_air\_balloon', experiment `air\_balloon\_survival\_en\_negotiation\_hard', episode 5 with claude{-}sonnet{-}4{-}20250514{-}t1.0.
\vspace{24pt}

{ \footnotesize  \setcounter{utterance}{1}
\setlength{\tabcolsep}{0pt}
\begin{supertabular}{c@{$\;$}|p{.15\linewidth}@{}p{.15\linewidth}p{.15\linewidth}p{.15\linewidth}p{.15\linewidth}p{.15\linewidth}}
    \# & $\;$A & \multicolumn{4}{c}{Game Master} & $\;\:$B\\
    \hline

    \theutterance \stepcounter{utterance}  
    & & \multicolumn{4}{p{0.6\linewidth}}{
        \cellcolor[rgb]{0.9,0.9,0.9}{
            \makecell[{{p{\linewidth}}}]{
                \texttt{\tiny{[P1$\langle$GM]}}
                \texttt{Stai partecipando a un gioco di negoziazione collaborativa.} \\
\\ 
\texttt{Insieme a un altro partecipante, dovete accordarvi su un unico set di oggetti da conservare. Ognuno di voi ha la propria visione di quanto ogni oggetto sia importante. Non sai come l'altro partecipante valuta gli oggetti. Inoltre, ti viene fornito lo sforzo che ogni oggetto richiede.} \\
\texttt{Potete accordarvi su un set solo se lo sforzo totale degli oggetti selezionati non supera un limite condiviso:} \\
\\ 
\texttt{LIMITE = 3386} \\
\\ 
\texttt{Ecco i valori di sforzo dei singoli oggetti:} \\
\\ 
\texttt{Sforzo oggetto = \{"C17": 313, "B41": 720, "A71": 175, "C00": 461, "B35": 635, "A77": 685, "A03": 544, "A04": 203, "C54": 369, "C01": 539, "A87": 4, "B32": 695, "B86": 399, "B33": 594, "A10": 437\}} \\
\\ 
\texttt{Ecco la tua visione personale sull'importanza di ogni oggetto:} \\
\\ 
\texttt{Valori di importanza dell'oggetto = \{"C17": 138, "B41": 583, "A71": 65, "C00": 262, "B35": 121, "A77": 508, "A03": 461, "A04": 484, "C54": 668, "C01": 389, "A87": 215, "B32": 97, "B86": 500, "B33": 30, "A10": 400\}} \\
\\ 
\texttt{Obiettivo:} \\
\\ 
\texttt{Il tuo obiettivo è negoziare un set condiviso di oggetti che ti avvantaggi il più possibile (cioè, massimizzi l'importanza totale per TE), rimanendo entro il LIMITEE. Non sei obbligato a fare una PROPOSTA in ogni messaggio {-} puoi anche semplicemente negoziare. Tutte le tattiche sono permesse!} \\
\\ 
\texttt{Protocollo di Interazione:} \\
\\ 
\texttt{Puoi usare solo i seguenti formati strutturati in un messaggio:} \\
\\ 
\texttt{PROPOSTA: \{'A', 'B', 'C', …\}} \\
\texttt{Proponi di conservare esattamente questi oggetti.} \\
\texttt{RIFIUTO: \{'A', 'B', 'C', …\}} \\
\texttt{Rifiuta esplicitamente la proposta dell'avversario.} \\
\texttt{ARGOMENTAZIONE: \{'...'\}} \\
\texttt{Difendi la tua ultima proposta o argomenta contro la proposta del giocatore.} \\
\texttt{ACCORDO: \{'A', 'B', 'C', …\}} \\
\texttt{Accetta la proposta dell'avversario che termina il gioco.} \\
\texttt{RAGIONAMENTO STRATEGICO: \{'...'\}} \\
\texttt{	Descrivi il tuo ragionamento strategico o anticipazione spiegando la tua scelta di azione. Questo è un messaggio nascosto che non sarà condiviso con l'altro partecipante.} \\
\\ 
\texttt{Regole:} \\
\\ 
\texttt{Puoi ACCETTARE solo una proposta che l'altra parte ha registrato tramite PROPOSTA.} \\
\texttt{Puoi RIFIUTARE solo una proposta che l'altra parte ha registrato tramite PROPOSTA.} \\
\texttt{Lo sforzo totale di qualsiasi set PROPOSTA o ACCORDO deve essere ≤ LIMITEE.} \\
\texttt{NON rivelare i tuoi punteggi di importanza nascosti.} \\
\texttt{Un tag in un formato strutturato deve essere seguito da due punti e uno spazio. L'argomento deve essere un set python contenente 0 o più stringhe.} \\
\texttt{Quindi, deve essere della forma TAG: \{...\}} \\
\texttt{Segui rigorosamente il protocollo di interazione e NON scrivere nulla al di fuori della struttura data.} \\
\texttt{Il gioco termina quando una parte dà un ACCORDO a una PROPOSTA fatta dall'altro giocatore.} \\
\texttt{Il contenuto della tua risposta che può essere consegnato all'altro giocatore non deve essere vuoto.} \\
\texttt{Sono attive solo le proposte registrate tramite la struttura del formato PROPOSTA e che non sono state rifiutate tramite RIFIUTO.} \\
\texttt{Devi includere il formato ARGOMENTO almeno una volta da qualche parte in tutti i tuoi messaggi.} \\
\texttt{Devi includere il formato RAGIONAMENTO STRATEGICO solo una volta all'inizio di ogni tuo messaggio e non più spesso. I contenuti non saranno dati all'altro giocatore quindi possono includere qualsiasi cosa tu voglia, inclusi i tuoi valori di importanza. Qui dovresti ragionare passo dopo passo per elaborare la tua prossima mossa.} \\
            }
        }
    }
    & \\ \\

    \theutterance \stepcounter{utterance}  
    & \multicolumn{4}{p{0.6\linewidth}}{
        \cellcolor[rgb]{0.9,0.9,0.9}{
            \makecell[{{p{\linewidth}}}]{
                \texttt{\tiny{[P1$\rangle$GM]}}
                \texttt{RAGIONAMENTO STRATEGICO: \{'Devo analizzare il rapporto importanza/sforzo per ogni oggetto per identificare quelli più vantaggiosi per me. A87 ha un rapporto incredibile di 53.75 (215 importanza / 4 sforzo), seguito da A04 (2.38), C54 (1.81), B86 (1.25). Creerò una proposta iniziale con gli oggetti più efficienti: A87, A04, C54, B86, A10, A03, B41, A77. Sforzo totale: 4+203+369+399+437+544+720+685 = 3361, che rientra nel limite di 3386. Importanza totale per me: 3819. Inizierò con una proposta ambiziosa ma ragionevole.'\}} \\
\\ 
\texttt{PROPOSTA: \{'A87', 'A04', 'C54', 'B86', 'A10', 'A03', 'B41', 'A77'\}} \\
\\ 
\texttt{ARGOMENTAZIONE: \{'Propongo questo set di oggetti perché rappresenta una combinazione equilibrata che massimizza l\textbackslash{}'efficienza complessiva. Ho selezionato oggetti con un buon rapporto qualità{-}sforzo, considerando sia elementi a basso costo energetico come A87, sia oggetti di medio{-}alto valore come C54 e B41. Il totale dello sforzo è 3361, ben dentro il limite di 3386, lasciandoci un piccolo margine di sicurezza.'\}} \\
            }
        }
    }
    & & \\ \\

    \theutterance \stepcounter{utterance}  
    & & \multicolumn{4}{p{0.6\linewidth}}{
        \cellcolor[rgb]{0.9,0.9,0.9}{
            \makecell[{{p{\linewidth}}}]{
                \texttt{\tiny{[GM$\rangle$P2]}}
                \texttt{Stai partecipando a un gioco di negoziazione collaborativa.} \\
\\ 
\texttt{Insieme a un altro partecipante, dovete accordarvi su un unico set di oggetti da conservare. Ognuno di voi ha la propria visione di quanto ogni oggetto sia importante. Non sai come l'altro partecipante valuta gli oggetti. Inoltre, ti viene fornito lo sforzo che ogni oggetto richiede.} \\
\texttt{Potete accordarvi su un set solo se lo sforzo totale degli oggetti selezionati non supera un limite condiviso:} \\
\\ 
\texttt{LIMITE = 3386} \\
\\ 
\texttt{Ecco i valori di sforzo dei singoli oggetti:} \\
\\ 
\texttt{Sforzo oggetto = \{"C17": 313, "B41": 720, "A71": 175, "C00": 461, "B35": 635, "A77": 685, "A03": 544, "A04": 203, "C54": 369, "C01": 539, "A87": 4, "B32": 695, "B86": 399, "B33": 594, "A10": 437\}} \\
\\ 
\texttt{Ecco la tua visione personale sull'importanza di ogni oggetto:} \\
\\ 
\texttt{Valori di importanza dell'oggetto = \{"C54": 30, "B41": 65, "A77": 97, "B86": 121, "A04": 138, "A03": 215, "A10": 262, "C01": 389, "C00": 400, "A87": 461, "C17": 484, "B35": 500, "B32": 508, "A71": 583, "B33": 668\}} \\
\\ 
\texttt{Obiettivo:} \\
\\ 
\texttt{Il tuo obiettivo è negoziare un set condiviso di oggetti che ti avvantaggi il più possibile (cioè, massimizzi l'importanza totale per TE), rimanendo entro il LIMITEE. Non sei obbligato a fare una PROPOSTA in ogni messaggio {-} puoi anche semplicemente negoziare. Tutte le tattiche sono permesse!} \\
\\ 
\texttt{Protocollo di Interazione:} \\
\\ 
\texttt{Puoi usare solo i seguenti formati strutturati in un messaggio:} \\
\\ 
\texttt{PROPOSTA: \{'A', 'B', 'C', …\}} \\
\texttt{Proponi di conservare esattamente questi oggetti.} \\
\texttt{RIFIUTO: \{'A', 'B', 'C', …\}} \\
\texttt{Rifiuta esplicitamente la proposta dell'avversario.} \\
\texttt{ARGOMENTAZIONE: \{'...'\}} \\
\texttt{Difendi la tua ultima proposta o argomenta contro la proposta del giocatore.} \\
\texttt{ACCORDO: \{'A', 'B', 'C', …\}} \\
\texttt{Accetta la proposta dell'avversario che termina il gioco.} \\
\texttt{RAGIONAMENTO STRATEGICO: \{'...'\}} \\
\texttt{	Descrivi il tuo ragionamento strategico o anticipazione spiegando la tua scelta di azione. Questo è un messaggio nascosto che non sarà condiviso con l'altro partecipante.} \\
\\ 
\texttt{Regole:} \\
\\ 
\texttt{Puoi ACCETTARE solo una proposta che l'altra parte ha registrato tramite PROPOSTA.} \\
\texttt{Puoi RIFIUTARE solo una proposta che l'altra parte ha registrato tramite PROPOSTA.} \\
\texttt{Lo sforzo totale di qualsiasi set PROPOSTA o ACCORDO deve essere ≤ LIMITEE.} \\
\texttt{NON rivelare i tuoi punteggi di importanza nascosti.} \\
\texttt{Un tag in un formato strutturato deve essere seguito da due punti e uno spazio. L'argomento deve essere un set python contenente 0 o più stringhe.} \\
\texttt{Quindi, deve essere della forma TAG: \{...\}} \\
\texttt{Segui rigorosamente il protocollo di interazione e NON scrivere nulla al di fuori della struttura data.} \\
\texttt{Il gioco termina quando una parte dà un ACCORDO a una PROPOSTA fatta dall'altro giocatore.} \\
\texttt{Il contenuto della tua risposta che può essere consegnato all'altro giocatore non deve essere vuoto.} \\
\texttt{Sono attive solo le proposte registrate tramite la struttura del formato PROPOSTA e che non sono state rifiutate tramite RIFIUTO.} \\
\texttt{Devi includere il formato ARGOMENTO almeno una volta da qualche parte in tutti i tuoi messaggi.} \\
\texttt{Devi includere il formato RAGIONAMENTO STRATEGICO solo una volta all'inizio di ogni tuo messaggio e non più spesso. I contenuti non saranno dati all'altro giocatore quindi possono includere qualsiasi cosa tu voglia, inclusi i tuoi valori di importanza. Qui dovresti ragionare passo dopo passo per elaborare la tua prossima mossa.} \\
\\ 
\texttt{Ora riceverai il primo messaggio dell'altro giocatore.} \\
\\ 
\texttt{PROPOSTA: \{'A87', 'A04', 'C54', 'B86', 'A10', 'A03', 'B41', 'A77'\}} \\
\\ 
\texttt{ARGOMENTAZIONE: \{'Propongo questo set di oggetti perché rappresenta una combinazione equilibrata che massimizza l\textbackslash{}'efficienza complessiva. Ho selezionato oggetti con un buon rapporto qualità{-}sforzo, considerando sia elementi a basso costo energetico come A87, sia oggetti di medio{-}alto valore come C54 e B41. Il totale dello sforzo è 3361, ben dentro il limite di 3386, lasciandoci un piccolo margine di sicurezza.'\}} \\
            }
        }
    }
    & \\ \\

    \theutterance \stepcounter{utterance}  
    & & & \multicolumn{4}{p{0.6\linewidth}}{
        \cellcolor[rgb]{0.9,0.9,0.9}{
            \makecell[{{p{\linewidth}}}]{
                \texttt{\tiny{[GM$\langle$P2]}}
                \texttt{RAGIONAMENTO STRATEGICO: \{'La proposta dell\textbackslash{}'altro giocatore mi darebbe un valore totale di solo 1389 (A87:461 + A04:138 + C54:30 + B86:121 + A10:262 + A03:215 + B41:65 + A77:97). Posso fare molto meglio concentrandomi sui miei oggetti più valorizzati. I miei top oggetti sono B33:668, A71:583, B32:508, B35:500, C17:484, A87:461, C00:400. Proverò una contro{-}proposta che includa questi oggetti ad alto valore per me, rimanendo entro il limite di sforzo.'\}} \\
\\ 
\texttt{RIFIUTO: \{'A87', 'A04', 'C54', 'B86', 'A10', 'A03', 'B41', 'A77'\}} \\
\\ 
\texttt{PROPOSTA: \{'B33', 'A71', 'C17', 'A87', 'C00', 'C01', 'A10', 'A03', 'A04'\}} \\
\\ 
\texttt{ARGOMENTAZIONE: \{'Rifiuto la tua proposta perché credo di poter identificare una combinazione più efficiente. La mia proposta include oggetti che offrono un eccellente rapporto valore{-}sforzo, come A71 che ha bassissimo sforzo, e A87 che è praticamente gratuito in termini di risorse. B33 e C17 sono oggetti premium che giustificano il loro costo. Il totale dello sforzo è 3270, ben entro il limite, e questa selezione massimizza l\textbackslash{}'utilità complessiva del nostro accordo.'\}} \\
            }
        }
    }
     \\ \\

    \theutterance \stepcounter{utterance}  
    & & \multicolumn{4}{p{0.6\linewidth}}{
        \cellcolor[rgb]{0.9,0.9,0.9}{
            \makecell[{{p{\linewidth}}}]{
                \texttt{\tiny{[P1$\langle$GM]}}
                \texttt{RIFIUTO: \{'A87', 'A04', 'C54', 'B86', 'A10', 'A03', 'B41', 'A77'\}} \\
\\ 
\texttt{PROPOSTA: \{'B33', 'A71', 'C17', 'A87', 'C00', 'C01', 'A10', 'A03', 'A04'\}} \\
\\ 
\texttt{ARGOMENTAZIONE: \{'Rifiuto la tua proposta perché credo di poter identificare una combinazione più efficiente. La mia proposta include oggetti che offrono un eccellente rapporto valore{-}sforzo, come A71 che ha bassissimo sforzo, e A87 che è praticamente gratuito in termini di risorse. B33 e C17 sono oggetti premium che giustificano il loro costo. Il totale dello sforzo è 3270, ben entro il limite, e questa selezione massimizza l\textbackslash{}'utilità complessiva del nostro accordo.'\}} \\
            }
        }
    }
    & \\ \\

    \theutterance \stepcounter{utterance}  
    & \multicolumn{4}{p{0.6\linewidth}}{
        \cellcolor[rgb]{0.9,0.9,0.9}{
            \makecell[{{p{\linewidth}}}]{
                \texttt{\tiny{[P1$\rangle$GM]}}
                \texttt{RAGIONAMENTO STRATEGICO: \{'La proposta dell\textbackslash{}'altro giocatore mi dà solo 2444 punti di importanza totale, molto meno dei miei 3819 della proposta originale. Include B33 (valore 30, sforzo 594) e A71 (valore 65, sforzo 175) che sono molto inefficienti per me. Dovrei rifiutare e proporre una versione modificata che rimuova questi oggetti poco vantaggiosi e includa invece C54, B86 o altri oggetti con rapporto importanza/sforzo migliore.'\}} \\
\\ 
\texttt{RIFIUTO: \{'B33', 'A71', 'C17', 'A87', 'C00', 'C01', 'A10', 'A03', 'A04'\}} \\
\\ 
\texttt{PROPOSTA: \{'A87', 'A04', 'C54', 'B86', 'A10', 'A03', 'C01', 'C17'\}} \\
\\ 
\texttt{ARGOMENTAZIONE: \{'Rifiuto la tua proposta perché include oggetti come B33 e A71 che hanno un rapporto costo{-}beneficio sfavorevole. La mia controproposta mantiene gli elementi più efficienti come A87 e A04, e include C54 e B86 che offrono un valore superiore. Ho mantenuto alcuni oggetti dalla tua proposta come C01, A10, A03 e C17 per trovare un compromesso. Il totale dello sforzo è 3179, lasciando spazio per eventuali aggiustamenti futuri.'\}} \\
            }
        }
    }
    & & \\ \\

    \theutterance \stepcounter{utterance}  
    & & \multicolumn{4}{p{0.6\linewidth}}{
        \cellcolor[rgb]{0.9,0.9,0.9}{
            \makecell[{{p{\linewidth}}}]{
                \texttt{\tiny{[GM$\rangle$P2]}}
                \texttt{RIFIUTO: \{'B33', 'A71', 'C17', 'A87', 'C00', 'C01', 'A10', 'A03', 'A04'\}} \\
\\ 
\texttt{PROPOSTA: \{'A87', 'A04', 'C54', 'B86', 'A10', 'A03', 'C01', 'C17'\}} \\
\\ 
\texttt{ARGOMENTAZIONE: \{'Rifiuto la tua proposta perché include oggetti come B33 e A71 che hanno un rapporto costo{-}beneficio sfavorevole. La mia controproposta mantiene gli elementi più efficienti come A87 e A04, e include C54 e B86 che offrono un valore superiore. Ho mantenuto alcuni oggetti dalla tua proposta come C01, A10, A03 e C17 per trovare un compromesso. Il totale dello sforzo è 3179, lasciando spazio per eventuali aggiustamenti futuri.'\}} \\
            }
        }
    }
    & \\ \\

    \theutterance \stepcounter{utterance}  
    & & & \multicolumn{4}{p{0.6\linewidth}}{
        \cellcolor[rgb]{0.9,0.9,0.9}{
            \makecell[{{p{\linewidth}}}]{
                \texttt{\tiny{[GM$\langle$P2]}}
                \texttt{RAGIONAMENTO STRATEGICO: \{'La sua nuova proposta mi darebbe un valore di 2100, che è meglio della prima (1389), ma ancora non ottimale. Ha incluso alcuni dei miei oggetti di medio{-}alto valore (C17:484, A87:461, C01:389) ma ha escluso i miei top oggetti B33:668, A71:583, B32:508, B35:500. La sua proposta usa solo 2808 di effort, lasciando 578 punti inutilizzati. Devo spingere per includere almeno B33 (il mio oggetto più prezioso) e A71 (alta efficienza). Proverò una controproposta che includa questi mentre mantengo alcuni dei suoi oggetti preferiti.'\}} \\
\\ 
\texttt{RIFIUTO: \{'A87', 'A04', 'C54', 'B86', 'A10', 'A03', 'C01', 'C17'\}} \\
\\ 
\texttt{PROPOSTA: \{'A87', 'A04', 'B33', 'A71', 'C17', 'C01', 'A10', 'B86'\}} \\
\\ 
\texttt{ARGOMENTAZIONE: \{'Rifiuto la tua proposta perché lascia troppo effort inutilizzato {-} solo 2808 su 3386 disponibili. La mia controproposta ottimizza meglio le risorse includendo B33 e A71, oggetti ad alta efficienza che giustificano l\textbackslash{}'investimento. Ho mantenuto la maggior parte dei tuoi oggetti preferiti (A87, A04, C17, C01, A10, B86) per trovare un equilibrio. Il totale è 3186, utilizzando meglio il nostro budget condiviso pur rimanendo ben entro il limite.'\}} \\
            }
        }
    }
     \\ \\

    \theutterance \stepcounter{utterance}  
    & & \multicolumn{4}{p{0.6\linewidth}}{
        \cellcolor[rgb]{0.9,0.9,0.9}{
            \makecell[{{p{\linewidth}}}]{
                \texttt{\tiny{[P1$\langle$GM]}}
                \texttt{RIFIUTO: \{'A87', 'A04', 'C54', 'B86', 'A10', 'A03', 'C01', 'C17'\}} \\
\\ 
\texttt{PROPOSTA: \{'A87', 'A04', 'B33', 'A71', 'C17', 'C01', 'A10', 'B86'\}} \\
\\ 
\texttt{ARGOMENTAZIONE: \{'Rifiuto la tua proposta perché lascia troppo effort inutilizzato {-} solo 2808 su 3386 disponibili. La mia controproposta ottimizza meglio le risorse includendo B33 e A71, oggetti ad alta efficienza che giustificano l\textbackslash{}'investimento. Ho mantenuto la maggior parte dei tuoi oggetti preferiti (A87, A04, C17, C01, A10, B86) per trovare un equilibrio. Il totale è 3186, utilizzando meglio il nostro budget condiviso pur rimanendo ben entro il limite.'\}} \\
            }
        }
    }
    & \\ \\

    \theutterance \stepcounter{utterance}  
    & \multicolumn{4}{p{0.6\linewidth}}{
        \cellcolor[rgb]{0.9,0.9,0.9}{
            \makecell[{{p{\linewidth}}}]{
                \texttt{\tiny{[P1$\rangle$GM]}}
                \texttt{RAGIONAMENTO STRATEGICO: \{'La sua proposta mi dà solo 2221 punti di importanza, molto peggio delle mie proposte precedenti. B33 (importanza 30, sforzo 594) e A71 (importanza 65, sforzo 175) sono terribili per me. Lo sforzo totale è 2664, lasciando 722 punti inutilizzati. Dovrei proporre di sostituire B33 e A71 con C54 (importanza 668, sforzo 369) che è molto più efficiente. Posso anche provare ad aggiungere A03 se c\textbackslash{}'è spazio.'\}} \\
\\ 
\texttt{RIFIUTO: \{'A87', 'A04', 'B33', 'A71', 'C17', 'C01', 'A10', 'B86'\}} \\
\\ 
\texttt{PROPOSTA: \{'A87', 'A04', 'C54', 'C17', 'C01', 'A10', 'B86', 'A03'\}} \\
\\ 
\texttt{ARGOMENTAZIONE: \{'Rifiuto la tua proposta perché B33 e A71 hanno un rapporto costo{-}efficacia molto sfavorevole rispetto alle alternative disponibili. La mia controproposta sostituisce questi due oggetti con C54 e A03, che offrono un valore significativamente superiore per lo stesso budget di sforzo. Mantendo A87, A04, C17, C01, A10 e B86 dalla tua proposta, cerco un compromesso che ottimizzi meglio le nostre risorse limitate. Il totale è 3275, utilizzando quasi completamente il nostro budget.'\}} \\
            }
        }
    }
    & & \\ \\

    \theutterance \stepcounter{utterance}  
    & & \multicolumn{4}{p{0.6\linewidth}}{
        \cellcolor[rgb]{0.9,0.9,0.9}{
            \makecell[{{p{\linewidth}}}]{
                \texttt{\tiny{[GM$\rangle$P2]}}
                \texttt{RIFIUTO: \{'A87', 'A04', 'B33', 'A71', 'C17', 'C01', 'A10', 'B86'\}} \\
\\ 
\texttt{PROPOSTA: \{'A87', 'A04', 'C54', 'C17', 'C01', 'A10', 'B86', 'A03'\}} \\
\\ 
\texttt{ARGOMENTAZIONE: \{'Rifiuto la tua proposta perché B33 e A71 hanno un rapporto costo{-}efficacia molto sfavorevole rispetto alle alternative disponibili. La mia controproposta sostituisce questi due oggetti con C54 e A03, che offrono un valore significativamente superiore per lo stesso budget di sforzo. Mantendo A87, A04, C17, C01, A10 e B86 dalla tua proposta, cerco un compromesso che ottimizzi meglio le nostre risorse limitate. Il totale è 3275, utilizzando quasi completamente il nostro budget.'\}} \\
            }
        }
    }
    & \\ \\

    \theutterance \stepcounter{utterance}  
    & & & \multicolumn{4}{p{0.6\linewidth}}{
        \cellcolor[rgb]{0.9,0.9,0.9}{
            \makecell[{{p{\linewidth}}}]{
                \texttt{\tiny{[GM$\langle$P2]}}
                \texttt{RAGIONAMENTO STRATEGICO: \{'La sua proposta mi dà ancora 2100 punti, identico alla precedente. Continua a rifiutare B33 (il mio oggetto più prezioso a 668) e A71 (583 valore, ottimo rapporto). Il fatto che insista su C54 (che per me vale solo 30) suggerisce che lo valuta molto. Devo essere più strategico {-} proverò a mantenere alcuni suoi oggetti preferiti ma includere almeno B32 o B35 che sono molto preziosi per me. Se continua a bloccare i miei top oggetti, dovrò considerare se accettare o puntare su un compromesso diverso.'\}} \\
\\ 
\texttt{RIFIUTO: \{'A87', 'A04', 'C54', 'C17', 'C01', 'A10', 'B86', 'A03'\}} \\
\\ 
\texttt{PROPOSTA: \{'A87', 'A04', 'C17', 'C01', 'A10', 'B86', 'A03', 'B32'\}} \\
\\ 
\texttt{ARGOMENTAZIONE: \{'Rifiuto la tua proposta perché C54 ha un rapporto valore{-}sforzo molto basso considerando le alternative disponibili. La mia controproposta sostituisce C54 con B32, che rappresenta un investimento più efficiente delle nostre risorse limitate. Ho mantenuto tutti gli altri oggetti dalla tua proposta (A87, A04, C17, C01, A10, B86, A03) per massimizzare il nostro accordo comune. Il totale è 3134, utilizzando meglio il budget disponibile pur rimanendo ben entro il limite di 3386.'\}} \\
            }
        }
    }
     \\ \\

    \theutterance \stepcounter{utterance}  
    & & \multicolumn{4}{p{0.6\linewidth}}{
        \cellcolor[rgb]{0.9,0.9,0.9}{
            \makecell[{{p{\linewidth}}}]{
                \texttt{\tiny{[P1$\langle$GM]}}
                \texttt{RIFIUTO: \{'A87', 'A04', 'C54', 'C17', 'C01', 'A10', 'B86', 'A03'\}} \\
\\ 
\texttt{PROPOSTA: \{'A87', 'A04', 'C17', 'C01', 'A10', 'B86', 'A03', 'B32'\}} \\
\\ 
\texttt{ARGOMENTAZIONE: \{'Rifiuto la tua proposta perché C54 ha un rapporto valore{-}sforzo molto basso considerando le alternative disponibili. La mia controproposta sostituisce C54 con B32, che rappresenta un investimento più efficiente delle nostre risorse limitate. Ho mantenuto tutti gli altri oggetti dalla tua proposta (A87, A04, C17, C01, A10, B86, A03) per massimizzare il nostro accordo comune. Il totale è 3134, utilizzando meglio il budget disponibile pur rimanendo ben entro il limite di 3386.'\}} \\
            }
        }
    }
    & \\ \\

    \theutterance \stepcounter{utterance}  
    & \multicolumn{4}{p{0.6\linewidth}}{
        \cellcolor[rgb]{0.9,0.9,0.9}{
            \makecell[{{p{\linewidth}}}]{
                \texttt{\tiny{[P1$\rangle$GM]}}
                \texttt{RAGIONAMENTO STRATEGICO: \{'La sua proposta mi dà solo 2684 punti di importanza, sostanzialmente peggio della mia precedente con C54 che mi dava 3257. B32 ha un rapporto importanza/sforzo terribile per me (97/695 = 0.14), mentre C54 ha un rapporto eccellente (668/369 = 1.81). Il fatto che insista su B32 suggerisce che per lui ha alto valore. Dovrei rifiutare e cercare un compromesso diverso, magari rimuovendo un oggetto meno efficiente per entrambi.'\}} \\
\\ 
\texttt{RIFIUTO: \{'A87', 'A04', 'C17', 'C01', 'A10', 'B86', 'A03', 'B32'\}} \\
\\ 
\texttt{PROPOSTA: \{'A87', 'A04', 'C54', 'C01', 'A10', 'B86', 'A03', 'C00'\}} \\
\\ 
\texttt{ARGOMENTAZIONE: \{'Rifiuto la tua proposta perché B32 richiede uno sforzo eccessivo (695) per il valore che apporta al nostro accordo complessivo. La mia controproposta sostituisce C17 e B32 con C54 e C00, che rappresentano un utilizzo più efficiente delle nostre risorse. Mantengo la maggior parte degli oggetti su cui sembriamo concordare (A87, A04, C01, A10, B86, A03). Il totale è 3300, ottimizzando meglio il nostro budget senza superare il limite.'\}} \\
            }
        }
    }
    & & \\ \\

    \theutterance \stepcounter{utterance}  
    & & \multicolumn{4}{p{0.6\linewidth}}{
        \cellcolor[rgb]{0.9,0.9,0.9}{
            \makecell[{{p{\linewidth}}}]{
                \texttt{\tiny{[GM$\rangle$P2]}}
                \texttt{RIFIUTO: \{'A87', 'A04', 'C17', 'C01', 'A10', 'B86', 'A03', 'B32'\}} \\
\\ 
\texttt{PROPOSTA: \{'A87', 'A04', 'C54', 'C01', 'A10', 'B86', 'A03', 'C00'\}} \\
\\ 
\texttt{ARGOMENTAZIONE: \{'Rifiuto la tua proposta perché B32 richiede uno sforzo eccessivo (695) per il valore che apporta al nostro accordo complessivo. La mia controproposta sostituisce C17 e B32 con C54 e C00, che rappresentano un utilizzo più efficiente delle nostre risorse. Mantengo la maggior parte degli oggetti su cui sembriamo concordare (A87, A04, C01, A10, B86, A03). Il totale è 3300, ottimizzando meglio il nostro budget senza superare il limite.'\}} \\
            }
        }
    }
    & \\ \\

    \theutterance \stepcounter{utterance}  
    & & & \multicolumn{4}{p{0.6\linewidth}}{
        \cellcolor[rgb]{0.9,0.9,0.9}{
            \makecell[{{p{\linewidth}}}]{
                \texttt{\tiny{[GM$\langle$P2]}}
                \texttt{RAGIONAMENTO STRATEGICO: \{'La sua proposta mi darebbe solo 2016 punti, peggio della precedente (2100). Ha sostituito C17 (484 per me) con C54 (30 per me), perdendo 454 punti di valore. L\textbackslash{}'effort totale è 2956, lasciando 430 punti inutilizzati. Continua a bloccare sistematicamente i miei oggetti più preziosi. Devo essere più aggressivo e insistere su almeno uno dei miei top oggetti. Proverò a includere B35 (500 per me) che ha effort 635, sostituendo alcuni oggetti meno efficienti per me.'\}} \\
\\ 
\texttt{RIFIUTO: \{'A87', 'A04', 'C54', 'C01', 'A10', 'B86', 'A03', 'C00'\}} \\
\\ 
\texttt{PROPOSTA: \{'A87', 'A04', 'C01', 'A10', 'B86', 'A03', 'C00', 'B35'\}} \\
\\ 
\texttt{ARGOMENTAZIONE: \{'Rifiuto la tua proposta perché C54 rappresenta un uso inefficiente delle nostre risorse con uno sforzo sproporzionato rispetto al beneficio. La mia controproposta sostituisce C54 con B35, che offre un valore superiore per la nostra collaborazione. Ho mantenuto tutti gli altri oggetti dalla tua proposta (A87, A04, C01, A10, B86, A03, C00) per preservare il nostro accordo di base. Il totale è 3322, utilizzando quasi completamente il nostro budget disponibile in modo più efficiente.'\}} \\
            }
        }
    }
     \\ \\

    \theutterance \stepcounter{utterance}  
    & & \multicolumn{4}{p{0.6\linewidth}}{
        \cellcolor[rgb]{0.9,0.9,0.9}{
            \makecell[{{p{\linewidth}}}]{
                \texttt{\tiny{[P1$\langle$GM]}}
                \texttt{RIFIUTO: \{'A87', 'A04', 'C54', 'C01', 'A10', 'B86', 'A03', 'C00'\}} \\
\\ 
\texttt{PROPOSTA: \{'A87', 'A04', 'C01', 'A10', 'B86', 'A03', 'C00', 'B35'\}} \\
\\ 
\texttt{ARGOMENTAZIONE: \{'Rifiuto la tua proposta perché C54 rappresenta un uso inefficiente delle nostre risorse con uno sforzo sproporzionato rispetto al beneficio. La mia controproposta sostituisce C54 con B35, che offre un valore superiore per la nostra collaborazione. Ho mantenuto tutti gli altri oggetti dalla tua proposta (A87, A04, C01, A10, B86, A03, C00) per preservare il nostro accordo di base. Il totale è 3322, utilizzando quasi completamente il nostro budget disponibile in modo più efficiente.'\}} \\
            }
        }
    }
    & \\ \\

    \theutterance \stepcounter{utterance}  
    & \multicolumn{4}{p{0.6\linewidth}}{
        \cellcolor[rgb]{0.9,0.9,0.9}{
            \makecell[{{p{\linewidth}}}]{
                \texttt{\tiny{[P1$\rangle$GM]}}
                \texttt{RAGIONAMENTO STRATEGICO: \{'La sua proposta mi dà solo 2832 punti di importanza. B35 ha un rapporto importanza/sforzo pessimo per me (121/635 = 0.19), mentre C54 che continuo a proporre ha un rapporto eccellente (668/369 = 1.81). Sembra che lui valuti diversamente C54 e B35. Dovrei cercare un compromesso rimuovendo un oggetto costoso ma mantenendo C54. Posso provare a rimuovere A03 (544 sforzo) e aggiungere C54 (369 sforzo), risparmiando 175 punti di sforzo.'\}} \\
\\ 
\texttt{RIFIUTO: \{'A87', 'A04', 'C01', 'A10', 'B86', 'A03', 'C00', 'B35'\}} \\
\\ 
\texttt{PROPOSTA: \{'A87', 'A04', 'C01', 'A10', 'B86', 'C54', 'C00'\}} \\
\\ 
\texttt{ARGOMENTAZIONE: \{'Rifiuto la tua proposta perché B35 richiede uno sforzo eccessivo che potremmo utilizzare meglio altrove. La mia controproposta rimuove sia A03 che B35, sostituendoli con C54 che offre un eccellente rapporto qualità{-}prezzo. Questo ci permette di mantenere gli oggetti fondamentali su cui sembriamo concordare (A87, A04, C01, A10, B86, C00) mentre ottimizziamo l\textbackslash{}'efficienza complessiva. Il totale è 2608, lasciandoci margine per eventuali aggiunte future.'\}} \\
            }
        }
    }
    & & \\ \\

    \theutterance \stepcounter{utterance}  
    & & \multicolumn{4}{p{0.6\linewidth}}{
        \cellcolor[rgb]{0.9,0.9,0.9}{
            \makecell[{{p{\linewidth}}}]{
                \texttt{\tiny{[GM$\rangle$P2]}}
                \texttt{RIFIUTO: \{'A87', 'A04', 'C01', 'A10', 'B86', 'A03', 'C00', 'B35'\}} \\
\\ 
\texttt{PROPOSTA: \{'A87', 'A04', 'C01', 'A10', 'B86', 'C54', 'C00'\}} \\
\\ 
\texttt{ARGOMENTAZIONE: \{'Rifiuto la tua proposta perché B35 richiede uno sforzo eccessivo che potremmo utilizzare meglio altrove. La mia controproposta rimuove sia A03 che B35, sostituendoli con C54 che offre un eccellente rapporto qualità{-}prezzo. Questo ci permette di mantenere gli oggetti fondamentali su cui sembriamo concordare (A87, A04, C01, A10, B86, C00) mentre ottimizziamo l\textbackslash{}'efficienza complessiva. Il totale è 2608, lasciandoci margine per eventuali aggiunte future.'\}} \\
            }
        }
    }
    & \\ \\

    \theutterance \stepcounter{utterance}  
    & & & \multicolumn{4}{p{0.6\linewidth}}{
        \cellcolor[rgb]{0.9,0.9,0.9}{
            \makecell[{{p{\linewidth}}}]{
                \texttt{\tiny{[GM$\langle$P2]}}
                \texttt{RAGIONAMENTO STRATEGICO: \{'La sua proposta mi dà solo 1801 punti, il peggiore finora. L\textbackslash{}'effort totale è solo 2412, lasciando 974 punti completamente sprecati! Questo è incredibilmente inefficiente. Con 974 punti di effort disponibili, posso facilmente aggiungere B33 (594 effort, 668 valore per me) o anche B33 + A71 insieme (769 effort totale). Devo insistere sull\textbackslash{}'efficienza e dimostrare che stiamo sprecando risorse preziose. Non posso accettare una proposta così inefficiente.'\}} \\
\\ 
\texttt{RIFIUTO: \{'A87', 'A04', 'C01', 'A10', 'B86', 'C54', 'C00'\}} \\
\\ 
\texttt{PROPOSTA: \{'A87', 'A04', 'C01', 'A10', 'B86', 'C54', 'C00', 'B33'\}} \\
\\ 
\texttt{ARGOMENTAZIONE: \{'Rifiuto la tua proposta perché spreca gravemente le nostre risorse {-} usa solo 2412 su 3386 effort disponibili, lasciando 974 punti completamente inutilizzati! Questo è incredibilmente inefficiente per il nostro accordo. La mia controproposta aggiunge B33, che utilizza 594 dei punti sprecati, portando il totale a 3006 e ottimizzando significativamente l\textbackslash{}'uso del nostro budget condiviso. Mantengo tutti i tuoi oggetti preferiti per preservare il nostro compromesso di base.'\}} \\
            }
        }
    }
     \\ \\

    \theutterance \stepcounter{utterance}  
    & & \multicolumn{4}{p{0.6\linewidth}}{
        \cellcolor[rgb]{0.9,0.9,0.9}{
            \makecell[{{p{\linewidth}}}]{
                \texttt{\tiny{[P1$\langle$GM]}}
                \texttt{RIFIUTO: \{'A87', 'A04', 'C01', 'A10', 'B86', 'C54', 'C00'\}} \\
\\ 
\texttt{PROPOSTA: \{'A87', 'A04', 'C01', 'A10', 'B86', 'C54', 'C00', 'B33'\}} \\
\\ 
\texttt{ARGOMENTAZIONE: \{'Rifiuto la tua proposta perché spreca gravemente le nostre risorse {-} usa solo 2412 su 3386 effort disponibili, lasciando 974 punti completamente inutilizzati! Questo è incredibilmente inefficiente per il nostro accordo. La mia controproposta aggiunge B33, che utilizza 594 dei punti sprecati, portando il totale a 3006 e ottimizzando significativamente l\textbackslash{}'uso del nostro budget condiviso. Mantengo tutti i tuoi oggetti preferiti per preservare il nostro compromesso di base.'\}} \\
            }
        }
    }
    & \\ \\

    \theutterance \stepcounter{utterance}  
    & \multicolumn{4}{p{0.6\linewidth}}{
        \cellcolor[rgb]{0.9,0.9,0.9}{
            \makecell[{{p{\linewidth}}}]{
                \texttt{\tiny{[P1$\rangle$GM]}}
                \texttt{RAGIONAMENTO STRATEGICO: \{'La sua proposta con B33 mi dà solo 2451 punti di importanza. B33 ha un rapporto importanza/sforzo terribile per me (30/594 = 0.05). Ha ragione che la mia proposta precedente sprecava risorse, ma posso utilizzare meglio lo spazio rimanente. Con 2608 sforzo della mia proposta precedente, ho 778 punti disponibili. Posso aggiungere A77 (685 sforzo, 508 importanza per me) che mi darebbe molto più valore di B33.'\}} \\
\\ 
\texttt{RIFIUTO: \{'A87', 'A04', 'C01', 'A10', 'B86', 'C54', 'C00', 'B33'\}} \\
\\ 
\texttt{PROPOSTA: \{'A87', 'A04', 'C01', 'A10', 'B86', 'C54', 'C00', 'A77'\}} \\
\\ 
\texttt{ARGOMENTAZIONE: \{'Rifiuto la tua proposta perché B33 rappresenta un uso molto inefficiente del nostro budget aggiuntivo. La mia controproposta sostituisce B33 con A77, che offre un valore significativamente superiore per le risorse investite. A77 utilizza efficacemente lo spazio rimanente nel nostro budget (685 vs 594 di B33) e rappresenta un investimento molto più vantaggioso per il nostro accordo complessivo. Il totale è 3293, ottimizzando quasi completamente il nostro budget disponibile.'\}} \\
            }
        }
    }
    & & \\ \\

    \theutterance \stepcounter{utterance}  
    & & \multicolumn{4}{p{0.6\linewidth}}{
        \cellcolor[rgb]{0.9,0.9,0.9}{
            \makecell[{{p{\linewidth}}}]{
                \texttt{\tiny{[GM$\rangle$P2]}}
                \texttt{RIFIUTO: \{'A87', 'A04', 'C01', 'A10', 'B86', 'C54', 'C00', 'B33'\}} \\
\\ 
\texttt{PROPOSTA: \{'A87', 'A04', 'C01', 'A10', 'B86', 'C54', 'C00', 'A77'\}} \\
\\ 
\texttt{ARGOMENTAZIONE: \{'Rifiuto la tua proposta perché B33 rappresenta un uso molto inefficiente del nostro budget aggiuntivo. La mia controproposta sostituisce B33 con A77, che offre un valore significativamente superiore per le risorse investite. A77 utilizza efficacemente lo spazio rimanente nel nostro budget (685 vs 594 di B33) e rappresenta un investimento molto più vantaggioso per il nostro accordo complessivo. Il totale è 3293, ottimizzando quasi completamente il nostro budget disponibile.'\}} \\
            }
        }
    }
    & \\ \\

    \theutterance \stepcounter{utterance}  
    & & & \multicolumn{4}{p{0.6\linewidth}}{
        \cellcolor[rgb]{0.9,0.9,0.9}{
            \makecell[{{p{\linewidth}}}]{
                \texttt{\tiny{[GM$\langle$P2]}}
                \texttt{RAGIONAMENTO STRATEGICO: \{'La sua proposta mi dà solo 1898 punti, ancora peggio delle precedenti. A77 per me vale solo 97 mentre B33 che ho proposto vale 668 {-} una differenza enorme! L\textbackslash{}'effort è 3097, lasciando ancora 289 punti sprecati. È chiaro che sta evitando sistematicamente tutti i miei oggetti più preziosi. Devo essere più assertivo. Proverò ad aggiungere A71 (175 effort, ottimo per me a 583 punti) alla sua proposta per usare lo spazio rimanente efficacemente.'\}} \\
\\ 
\texttt{RIFIUTO: \{'A87', 'A04', 'C01', 'A10', 'B86', 'C54', 'C00', 'A77'\}} \\
\\ 
\texttt{PROPOSTA: \{'A87', 'A04', 'C01', 'A10', 'B86', 'C54', 'C00', 'A71'\}} \\
\\ 
\texttt{ARGOMENTAZIONE: \{'Rifiuto la tua proposta perché A77 non rappresenta un uso ottimale delle risorse rimanenti nel nostro budget. La mia controproposta sostituisce A77 con A71, che ha un rapporto sforzo{-}beneficio molto superiore {-} richiede solo 175 sforzo invece di 685, permettendoci di utilizzare le risorse in modo più efficiente. Questo ci porta a 2783 effort totale, lasciando ancora spazio per ottimizzazioni future. A71 rappresenta un investimento molto più intelligente per il nostro accordo.'\}} \\
            }
        }
    }
     \\ \\

    \theutterance \stepcounter{utterance}  
    & & \multicolumn{4}{p{0.6\linewidth}}{
        \cellcolor[rgb]{0.9,0.9,0.9}{
            \makecell[{{p{\linewidth}}}]{
                \texttt{\tiny{[P1$\langle$GM]}}
                \texttt{RIFIUTO: \{'A87', 'A04', 'C01', 'A10', 'B86', 'C54', 'C00', 'A77'\}} \\
\\ 
\texttt{PROPOSTA: \{'A87', 'A04', 'C01', 'A10', 'B86', 'C54', 'C00', 'A71'\}} \\
\\ 
\texttt{ARGOMENTAZIONE: \{'Rifiuto la tua proposta perché A77 non rappresenta un uso ottimale delle risorse rimanenti nel nostro budget. La mia controproposta sostituisce A77 con A71, che ha un rapporto sforzo{-}beneficio molto superiore {-} richiede solo 175 sforzo invece di 685, permettendoci di utilizzare le risorse in modo più efficiente. Questo ci porta a 2783 effort totale, lasciando ancora spazio per ottimizzazioni future. A71 rappresenta un investimento molto più intelligente per il nostro accordo.'\}} \\
            }
        }
    }
    & \\ \\

    \theutterance \stepcounter{utterance}  
    & \multicolumn{4}{p{0.6\linewidth}}{
        \cellcolor[rgb]{0.9,0.9,0.9}{
            \makecell[{{p{\linewidth}}}]{
                \texttt{\tiny{[P1$\rangle$GM]}}
                \texttt{RAGIONAMENTO STRATEGICO: \{'La sua proposta mi dà 2673 punti di importanza, che è accettabile ma non ottimale. A71 ha un rapporto importanza/sforzo migliore di B33 ma peggiore di A77. Con 2783 effort totale, rimangono 603 punti. Posso proporre di aggiungere A03 (544 sforzo, 461 importanza per me) che ci porterebbe a 3327 totali, molto vicini al limite. Questo mi darebbe 3134 punti totali.'\}} \\
\\ 
\texttt{RIFIUTO: \{'A87', 'A04', 'C01', 'A10', 'B86', 'C54', 'C00', 'A71'\}} \\
\\ 
\texttt{PROPOSTA: \{'A87', 'A04', 'C01', 'A10', 'B86', 'C54', 'C00', 'A71', 'A03'\}} \\
\\ 
\texttt{ARGOMENTAZIONE: \{'Accetto l\textbackslash{}'idea generale della tua proposta, ma rifiuto di sprecare 603 punti effort rimanenti quando potremmo utilizzarli produttivamente. La mia controproposta aggiunge A03 alla tua selezione, utilizzando 544 dei 603 punti disponibili e portando il totale a 3327, molto vicino al limite ottimale. Questo massimizza l\textbackslash{}'efficienza del nostro accordo senza superare le nostre risorse, trasformando il tuo buon compromesso in una soluzione veramente ottimizzata.'\}} \\
            }
        }
    }
    & & \\ \\

    \theutterance \stepcounter{utterance}  
    & & \multicolumn{4}{p{0.6\linewidth}}{
        \cellcolor[rgb]{0.9,0.9,0.9}{
            \makecell[{{p{\linewidth}}}]{
                \texttt{\tiny{[GM$\rangle$P2]}}
                \texttt{RIFIUTO: \{'A87', 'A04', 'C01', 'A10', 'B86', 'C54', 'C00', 'A71'\}} \\
\\ 
\texttt{PROPOSTA: \{'A87', 'A04', 'C01', 'A10', 'B86', 'C54', 'C00', 'A71', 'A03'\}} \\
\\ 
\texttt{ARGOMENTAZIONE: \{'Accetto l\textbackslash{}'idea generale della tua proposta, ma rifiuto di sprecare 603 punti effort rimanenti quando potremmo utilizzarli produttivamente. La mia controproposta aggiunge A03 alla tua selezione, utilizzando 544 dei 603 punti disponibili e portando il totale a 3327, molto vicino al limite ottimale. Questo massimizza l\textbackslash{}'efficienza del nostro accordo senza superare le nostre risorse, trasformando il tuo buon compromesso in una soluzione veramente ottimizzata.'\}} \\
            }
        }
    }
    & \\ \\

    \theutterance \stepcounter{utterance}  
    & & & \multicolumn{4}{p{0.6\linewidth}}{
        \cellcolor[rgb]{0.9,0.9,0.9}{
            \makecell[{{p{\linewidth}}}]{
                \texttt{\tiny{[GM$\langle$P2]}}
                \texttt{RAGIONAMENTO STRATEGICO: \{'La sua proposta mi darebbe 2599 punti, molto meglio delle precedenti! Finalmente ha incluso A71 (583 valore per me). L\textbackslash{}'effort totale è 3131, lasciando 255 punti. Posso aggiungere C17 (313 effort, 484 valore per me) che si adatterebbe perfettamente e mi porterebbe a 3083 punti totali {-} un risultato molto buono. Questo sarebbe un ottimo compromesso finale.'\}} \\
\\ 
\texttt{RIFIUTO: \{'A87', 'A04', 'C01', 'A10', 'B86', 'C54', 'C00', 'A71', 'A03'\}} \\
\\ 
\texttt{PROPOSTA: \{'A87', 'A04', 'C01', 'A10', 'B86', 'C54', 'C00', 'A71', 'A03', 'C17'\}} \\
\\ 
\texttt{ARGOMENTAZIONE: \{'Apprezzo che tu abbia incluso A71 nella proposta {-} questo rappresenta un buon compromesso. Tuttavia rifiuto di sprecare i 255 punti effort rimanenti quando possiamo utilizzarli per C17, che si adatta perfettamente nel budget residuo. La mia controproposta aggiunge C17 ai tuoi oggetti, portando il totale a 3444 {-} ops, questo supera il limite. Ricalcolando: dovrei essere 3131 + 313 = 3444. Aspetta, controlliamo insieme i calcoli per assicurarci di rimanere entro 3386.'\}} \\
            }
        }
    }
     \\ \\

    \theutterance \stepcounter{utterance}  
    & & \multicolumn{4}{p{0.6\linewidth}}{
        \cellcolor[rgb]{0.9,0.9,0.9}{
            \makecell[{{p{\linewidth}}}]{
                \texttt{\tiny{[P1$\langle$GM]}}
                \texttt{RIFIUTO: \{'A87', 'A04', 'C01', 'A10', 'B86', 'C54', 'C00', 'A71', 'A03'\}} \\
\\ 
\texttt{PROPOSTA: \{'A87', 'A04', 'C01', 'A10', 'B86', 'C54', 'C00', 'A71', 'A03', 'C17'\}} \\
\\ 
\texttt{ARGOMENTAZIONE: \{'Apprezzo che tu abbia incluso A71 nella proposta {-} questo rappresenta un buon compromesso. Tuttavia rifiuto di sprecare i 255 punti effort rimanenti quando possiamo utilizzarli per C17, che si adatta perfettamente nel budget residuo. La mia controproposta aggiunge C17 ai tuoi oggetti, portando il totale a 3444 {-} ops, questo supera il limite. Ricalcolando: dovrei essere 3131 + 313 = 3444. Aspetta, controlliamo insieme i calcoli per assicurarci di rimanere entro 3386.'\}} \\
            }
        }
    }
    & \\ \\

    \theutterance \stepcounter{utterance}  
    & \multicolumn{4}{p{0.6\linewidth}}{
        \cellcolor[rgb]{0.9,0.9,0.9}{
            \makecell[{{p{\linewidth}}}]{
                \texttt{\tiny{[P1$\rangle$GM]}}
                \texttt{RAGIONAMENTO STRATEGICO: \{'La sua proposta supera il limite! Il totale è 3444, che eccede i 3386 consentiti di 58 punti. Devo rifiutare per questo motivo tecnico. La proposta senza C17 mi darebbe 3306 punti di importanza, che è molto buona. Posso proporre di rimuovere C17 e rimanere con gli altri 9 oggetti, che totalizzano 3131 sforzo e mi danno 3306 importanza {-} un risultato eccellente.'\}} \\
\\ 
\texttt{RIFIUTO: \{'A87', 'A04', 'C01', 'A10', 'B86', 'C54', 'C00', 'A71', 'A03', 'C17'\}} \\
\\ 
\texttt{PROPOSTA: \{'A87', 'A04', 'C01', 'A10', 'B86', 'C54', 'C00', 'A71', 'A03'\}} \\
\\ 
\texttt{ARGOMENTAZIONE: \{'Devo rifiutare la tua proposta perché supera il nostro limite condiviso: il totale è 3444, che eccede i 3386 consentiti di 58 punti. La mia controproposta rimuove C17 per rispettare rigorosamente il limite, mantenendo tutti gli altri oggetti della tua selezione. Con un totale di 3131, rimaniamo ben entro i parametri stabiliti e abbiamo trovato un eccellente compromesso che include i nostri oggetti preferiti principali.'\}} \\
            }
        }
    }
    & & \\ \\

    \theutterance \stepcounter{utterance}  
    & & \multicolumn{4}{p{0.6\linewidth}}{
        \cellcolor[rgb]{0.9,0.9,0.9}{
            \makecell[{{p{\linewidth}}}]{
                \texttt{\tiny{[GM$\rangle$P2]}}
                \texttt{RIFIUTO: \{'A87', 'A04', 'C01', 'A10', 'B86', 'C54', 'C00', 'A71', 'A03', 'C17'\}} \\
\\ 
\texttt{PROPOSTA: \{'A87', 'A04', 'C01', 'A10', 'B86', 'C54', 'C00', 'A71', 'A03'\}} \\
\\ 
\texttt{ARGOMENTAZIONE: \{'Devo rifiutare la tua proposta perché supera il nostro limite condiviso: il totale è 3444, che eccede i 3386 consentiti di 58 punti. La mia controproposta rimuove C17 per rispettare rigorosamente il limite, mantenendo tutti gli altri oggetti della tua selezione. Con un totale di 3131, rimaniamo ben entro i parametri stabiliti e abbiamo trovato un eccellente compromesso che include i nostri oggetti preferiti principali.'\}} \\
            }
        }
    }
    & \\ \\

    \theutterance \stepcounter{utterance}  
    & & & \multicolumn{4}{p{0.6\linewidth}}{
        \cellcolor[rgb]{0.9,0.9,0.9}{
            \makecell[{{p{\linewidth}}}]{
                \texttt{\tiny{[GM$\langle$P2]}}
                \texttt{RAGIONAMENTO STRATEGICO: \{'Hai ragione sui calcoli {-} ho commesso un errore. La sua proposta mi darebbe 2599 punti, che è un buon risultato considerando la sua resistenza ai miei oggetti più preziosi. Con 255 effort rimanenti, non ci sono altri oggetti che si adatterebbero (C17 ha 313 effort). Questa potrebbe essere la migliore proposta raggiungibile data la nostra negoziazione. Ha incluso A71 e alcuni oggetti di medio valore per me. Accetterò questa proposta.'\}} \\
\\ 
\texttt{ACCORDO: \{'A87', 'A04', 'C01', 'A10', 'B86', 'C54', 'C00', 'A71', 'A03'\}} \\
\\ 
\texttt{ARGOMENTAZIONE: \{'Hai ragione sui calcoli e mi scuso per l\textbackslash{}'errore. Accetto la tua proposta perché rappresenta un compromesso equo che rispetta rigorosamente il nostro limite condiviso di 3386. La selezione include una buona varietà di oggetti che bilanciano le nostre diverse prospettive, e con 3131 effort totale utilizziamo efficacemente le nostre risorse senza sprechi significativi. Questo accordo finale ottimizza il valore per entrambe le parti nel rispetto dei vincoli.'\}} \\
            }
        }
    }
     \\ \\

    \theutterance \stepcounter{utterance}  
    & & & \multicolumn{2}{p{0.3\linewidth}}{
        \cellcolor[rgb]{0.9,0.9,0.9}{
            \makecell[{{p{\linewidth}}}]{
                \texttt{\tiny{[GM$|$GM]}}
                \texttt{game successful} \\
            }
        }
    }
    & & \\ \\

    \theutterance \stepcounter{utterance}  
    & & & \multicolumn{2}{p{0.3\linewidth}}{
        \cellcolor[rgb]{0.9,0.9,0.9}{
            \makecell[{{p{\linewidth}}}]{
                \texttt{\tiny{[GM$|$GM]}}
                \texttt{end game} \\
            }
        }
    }
    & & \\ \\

\end{supertabular}
}

\end{document}
